\hypertarget{appendix-1823a14330cb}{}\label{referencestart:appendix-1823a14330cb}
\section{Interval arithmetic}\label{appendix-1823a14330cb-interval-arithmetic}

\textbf{Curriculum depth:} Introduced · \textbf{First seen:} C004 · \textbf{Latest development:} C004

\begin{allmightroledefinition}

\subsection{Definition and mechanism}\label{appendix-1823a14330cb-definition-and-mechanism}

An interval encloses a possible exact value between lower and upper bounds. Correct outward rounding can preserve containment when the whole algorithm honors necessary assumptions.

\end{allmightroledefinition}

\subsection{Prerequisites}\label{appendix-1823a14330cb-prerequisites}

\begin{itemize}
\tightlist
\item
  \hyperref[appendix-637d164ae5ca]{Directed rounding modes} --- The signed rounding direction is state that a demonstration or enclosure algorithm must control.
\item
  \hyperref[appendix-f11d1a521253]{Floating-point environment and compiler modes} --- Compiler/runtime assumptions must preserve the intended direction and special-value semantics.
\end{itemize}

\subsection{Used by}\label{appendix-1823a14330cb-used-by}

Connect later dependent concepts here as their lessons arrive.

\subsection{Recall and next depth}\label{appendix-1823a14330cb-recall-and-next-depth}

Reconstruct the concrete example in the definition and identify its validity assumptions.

\textbf{Pending:} Directed rounding alone does not establish a valid numerical proof or robust predicate.

\subsection{Source and navigation}\label{appendix-1823a14330cb-source-and-navigation}

\hyperref[01-foundations-i-ch001-a-stretching-ruler-for-real-valued-quantities]{M004 teaching chapter} · \hyperref[appendix-fd090d982a15-m004---parts-61-76]{Per-part definitions} · \href{../07-Vision-XR-and-Playable-Worlds/main.pdf\#nameddest=appendix-4ad47bdd269a}{Dependency Map (Vision, XR, and Playable Worlds, Reference notes: Dependency Map)}
